\documentclass[12pt, a4paper]{article}

\usepackage[a4paper, margin=1in]{geometry}

\usepackage[utf8]{inputenc}
\usepackage[mathscr]{euscript}
\usepackage{amssymb,amsmath,amsthm,amsfonts}
\usepackage[shortlabels]{enumitem}
\usepackage{multicol,multirow}
\usepackage{lipsum}
\usepackage{balance}
\usepackage{calc}
\usepackage[colorlinks=true,citecolor=blue,linkcolor=blue]{hyperref}
\usepackage{import}
\usepackage{xifthen}
\usepackage{pdfpages}
\usepackage{transparent}
\usepackage{listings}

\newcommand{\incfig}[2][1.0]{
    \def\svgwidth{#1\columnwidth}
    \import{./figures/}{#2.pdf_tex}
}

\newcommand{\incimg}[2][1.0]{
  \includegraphics[width=#1\columnwidth]{./figures/#2}
}

\newlist{enumproof}{enumerate}{4}
\setlist[enumproof,1]{label=\arabic*., parsep=1em}
\setlist[enumproof,2]{label=\arabic{enumproofi}.\arabic*., parsep=1em}
\setlist[enumproof,3]{label=\arabic{enumproofi}.\arabic{enumproofii}.\arabic*., parsep=1em}
\setlist[enumproof,4]{label=\arabic{enumproofi}.\arabic{enumproofii}.\arabic{enumproofiii}.\arabic*., parsep=1em}

\renewcommand{\qedsymbol}{\ensuremath{\blacksquare}}

\lstdefinestyle{mystyle}{
  language=Python, % Set the language to Python
  commentstyle=\color{codegray}\itshape, % Python comments (#)
  keywordstyle=\color{orange}\bfseries,  % Python reserved keywords
  stringstyle=\color{mauve},             % Strings
  basicstyle={\ttfamily\footnotesize},   % Monospace font
  breakatwhitespace=true,         
  breaklines=true,                 
  captionpos=b,                    
  keepspaces=true,                 
  numbers=none,                    
  tabsize=2,
  showstringspaces=false,
  % Extra keyword groups for Python
  morekeywords=[2]{self, cls, __init__, __str__, __repr__}, 
  keywordstyle=[2]\color{identblue},     % Special identifiers / methods
  morekeywords=[3]{int, float, bool, str, list, dict, set, tuple, object, None, True, False},
  keywordstyle=[3]\color{codegreen},     % Built-in types
  morekeywords=[4]{print, input, open, range, len, type, isinstance, super, enumerate, zip, map, filter, sum, min, max, sorted},
  keywordstyle=[4]\color{cyan},          % Built-in functions
}

% Things Lie
\newcommand{\kb}{\mathfrak b}
\newcommand{\kg}{\mathfrak g}
\newcommand{\kh}{\mathfrak h}
\newcommand{\kn}{\mathfrak n}
\newcommand{\ku}{\mathfrak u}
\newcommand{\kz}{\mathfrak z}
\DeclareMathOperator{\Ext}{Ext} % Ext functor
\DeclareMathOperator{\Tor}{Tor} % Tor functor
\newcommand{\gl}{\opname{\mathfrak{gl}}} % frak gl group
\renewcommand{\sl}{\opname{\mathfrak{sl}}} % frak sl group chktex 6

% More script letters etc.
\newcommand{\SA}{\mathcal A}
\newcommand{\SB}{\mathcal B}
\newcommand{\SC}{\mathcal C}
\newcommand{\SF}{\mathcal F}
\newcommand{\SG}{\mathcal G}
\newcommand{\SH}{\mathcal H}
\newcommand{\OO}{\mathcal O}

\newcommand{\SCA}{\mathscr A}
\newcommand{\SCB}{\mathscr B}
\newcommand{\SCC}{\mathscr C}
\newcommand{\SCD}{\mathscr D}
\newcommand{\SCE}{\mathscr E}
\newcommand{\SCF}{\mathscr F}
\newcommand{\SCG}{\mathscr G}
\newcommand{\SCH}{\mathscr H}

% Mathfrak primes
\newcommand{\km}{\mathfrak m}
\newcommand{\kp}{\mathfrak p}
\newcommand{\kq}{\mathfrak q}

% number sets
\newcommand{\RR}[1][]{\ensuremath{\ifstrempty{#1}{\mathbb{R}}{\mathbb{R}^{#1}}}}
\newcommand{\NN}[1][]{\ensuremath{\ifstrempty{#1}{\mathbb{N}}{\mathbb{N}^{#1}}}}
\newcommand{\ZZ}[1][]{\ensuremath{\ifstrempty{#1}{\mathbb{Z}}{\mathbb{Z}^{#1}}}}
\newcommand{\QQ}[1][]{\ensuremath{\ifstrempty{#1}{\mathbb{Q}}{\mathbb{Q}^{#1}}}}
\newcommand{\CC}[1][]{\ensuremath{\ifstrempty{#1}{\mathbb{C}}{\mathbb{C}^{#1}}}}
\newcommand{\PP}[1][]{\ensuremath{\ifstrempty{#1}{\mathbb{P}}{\mathbb{P}^{#1}}}}
\newcommand{\HH}[1][]{\ensuremath{\ifstrempty{#1}{\mathbb{H}}{\mathbb{H}^{#1}}}}
\newcommand{\FF}[1][]{\ensuremath{\ifstrempty{#1}{\mathbb{F}}{\mathbb{F}^{#1}}}}
% expected value
\newcommand{\EE}{\ensuremath{\mathbb{E}}}
\newcommand{\charin}{\text{ char }}
\DeclareMathOperator{\sign}{sign}
\DeclareMathOperator{\Aut}{Aut}
\DeclareMathOperator{\Inn}{Inn}
\DeclareMathOperator{\Syl}{Syl}
\DeclareMathOperator{\Gal}{Gal}
\DeclareMathOperator{\GL}{GL} % General linear group
\DeclareMathOperator{\SL}{SL} % Special linear group

%---------------------------------------
% BlackBoard Math Fonts :-
%---------------------------------------

%Captital Letters
\newcommand{\bbA}{\mathbb{A}}	\newcommand{\bbB}{\mathbb{B}}
\newcommand{\bbC}{\mathbb{C}}	\newcommand{\bbD}{\mathbb{D}}
\newcommand{\bbE}{\mathbb{E}}	\newcommand{\bbF}{\mathbb{F}}
\newcommand{\bbG}{\mathbb{G}}	\newcommand{\bbH}{\mathbb{H}}
\newcommand{\bbI}{\mathbb{I}}	\newcommand{\bbJ}{\mathbb{J}}
\newcommand{\bbK}{\mathbb{K}}	\newcommand{\bbL}{\mathbb{L}}
\newcommand{\bbM}{\mathbb{M}}	\newcommand{\bbN}{\mathbb{N}}
\newcommand{\bbO}{\mathbb{O}}	\newcommand{\bbP}{\mathbb{P}}
\newcommand{\bbQ}{\mathbb{Q}}	\newcommand{\bbR}{\mathbb{R}}
\newcommand{\bbS}{\mathbb{S}}	\newcommand{\bbT}{\mathbb{T}}
\newcommand{\bbU}{\mathbb{U}}	\newcommand{\bbV}{\mathbb{V}}
\newcommand{\bbW}{\mathbb{W}}	\newcommand{\bbX}{\mathbb{X}}
\newcommand{\bbY}{\mathbb{Y}}	\newcommand{\bbZ}{\mathbb{Z}}

%---------------------------------------
% MathCal Fonts :-
%---------------------------------------

%Captital Letters
\newcommand{\mcA}{\mathcal{A}}	\newcommand{\mcB}{\mathcal{B}}
\newcommand{\mcC}{\mathcal{C}}	\newcommand{\mcD}{\mathcal{D}}
\newcommand{\mcE}{\mathcal{E}}	\newcommand{\mcF}{\mathcal{F}}
\newcommand{\mcG}{\mathcal{G}}	\newcommand{\mcH}{\mathcal{H}}
\newcommand{\mcI}{\mathcal{I}}	\newcommand{\mcJ}{\mathcal{J}}
\newcommand{\mcK}{\mathcal{K}}	\newcommand{\mcL}{\mathcal{L}}
\newcommand{\mcM}{\mathcal{M}}	\newcommand{\mcN}{\mathcal{N}}
\newcommand{\mcO}{\mathcal{O}}	\newcommand{\mcP}{\mathcal{P}}
\newcommand{\mcQ}{\mathcal{Q}}	\newcommand{\mcR}{\mathcal{R}}
\newcommand{\mcS}{\mathcal{S}}	\newcommand{\mcT}{\mathcal{T}}
\newcommand{\mcU}{\mathcal{U}}	\newcommand{\mcV}{\mathcal{V}}
\newcommand{\mcW}{\mathcal{W}}	\newcommand{\mcX}{\mathcal{X}}
\newcommand{\mcY}{\mathcal{Y}}	\newcommand{\mcZ}{\mathcal{Z}}


%---------------------------------------
% Bold Math Fonts :-
%---------------------------------------

%Captital Letters
\newcommand{\bmA}{\boldsymbol{A}}	\newcommand{\bmB}{\boldsymbol{B}}
\newcommand{\bmC}{\boldsymbol{C}}	\newcommand{\bmD}{\boldsymbol{D}}
\newcommand{\bmE}{\boldsymbol{E}}	\newcommand{\bmF}{\boldsymbol{F}}
\newcommand{\bmG}{\boldsymbol{G}}	\newcommand{\bmH}{\boldsymbol{H}}
\newcommand{\bmI}{\boldsymbol{I}}	\newcommand{\bmJ}{\boldsymbol{J}}
\newcommand{\bmK}{\boldsymbol{K}}	\newcommand{\bmL}{\boldsymbol{L}}
\newcommand{\bmM}{\boldsymbol{M}}	\newcommand{\bmN}{\boldsymbol{N}}
\newcommand{\bmO}{\boldsymbol{O}}	\newcommand{\bmP}{\boldsymbol{P}}
\newcommand{\bmQ}{\boldsymbol{Q}}	\newcommand{\bmR}{\boldsymbol{R}}
\newcommand{\bmS}{\boldsymbol{S}}	\newcommand{\bmT}{\boldsymbol{T}}
\newcommand{\bmU}{\boldsymbol{U}}	\newcommand{\bmV}{\boldsymbol{V}}
\newcommand{\bmW}{\boldsymbol{W}}	\newcommand{\bmX}{\boldsymbol{X}}
\newcommand{\bmY}{\boldsymbol{Y}}	\newcommand{\bmZ}{\boldsymbol{Z}}
%Small Letters
\newcommand{\bma}{\boldsymbol{a}}	\newcommand{\bmb}{\boldsymbol{b}}
\newcommand{\bmc}{\boldsymbol{c}}	\newcommand{\bmd}{\boldsymbol{d}}
\newcommand{\bme}{\boldsymbol{e}}	\newcommand{\bmf}{\boldsymbol{f}}
\newcommand{\bmg}{\boldsymbol{g}}	\newcommand{\bmh}{\boldsymbol{h}}
\newcommand{\bmi}{\boldsymbol{i}}	\newcommand{\bmj}{\boldsymbol{j}}
\newcommand{\bmk}{\boldsymbol{k}}	\newcommand{\bml}{\boldsymbol{l}}
\newcommand{\bmm}{\boldsymbol{m}}	\newcommand{\bmn}{\boldsymbol{n}}
\newcommand{\bmo}{\boldsymbol{o}}	\newcommand{\bmp}{\boldsymbol{p}}
\newcommand{\bmq}{\boldsymbol{q}}	\newcommand{\bmr}{\boldsymbol{r}}
\newcommand{\bms}{\boldsymbol{s}}	\newcommand{\bmt}{\boldsymbol{t}}
\newcommand{\bmu}{\boldsymbol{u}}	\newcommand{\bmv}{\boldsymbol{v}}
\newcommand{\bmw}{\boldsymbol{w}}	\newcommand{\bmx}{\boldsymbol{x}}
\newcommand{\bmy}{\boldsymbol{y}}	\newcommand{\bmz}{\boldsymbol{z}}

%---------------------------------------
% Scr Math Fonts :-
%---------------------------------------

\newcommand{\sA}{{\mathscr{A}}}   \newcommand{\sB}{{\mathscr{B}}}
\newcommand{\sC}{{\mathscr{C}}}   \newcommand{\sD}{{\mathscr{D}}}
\newcommand{\sE}{{\mathscr{E}}}   \newcommand{\sF}{{\mathscr{F}}}
\newcommand{\sG}{{\mathscr{G}}}   \newcommand{\sH}{{\mathscr{H}}}
\newcommand{\sI}{{\mathscr{I}}}   \newcommand{\sJ}{{\mathscr{J}}}
\newcommand{\sK}{{\mathscr{K}}}   \newcommand{\sL}{{\mathscr{L}}}
\newcommand{\sM}{{\mathscr{M}}}   \newcommand{\sN}{{\mathscr{N}}}
\newcommand{\sO}{{\mathscr{O}}}   \newcommand{\sP}{{\mathscr{P}}}
\newcommand{\sQ}{{\mathscr{Q}}}   \newcommand{\sR}{{\mathscr{R}}}
\newcommand{\sS}{{\mathscr{S}}}   \newcommand{\sT}{{\mathscr{T}}}
\newcommand{\sU}{{\mathscr{U}}}   \newcommand{\sV}{{\mathscr{V}}}
\newcommand{\sW}{{\mathscr{W}}}   \newcommand{\sX}{{\mathscr{X}}}
\newcommand{\sY}{{\mathscr{Y}}}   \newcommand{\sZ}{{\mathscr{Z}}}


%---------------------------------------
% Math Fraktur Font
%---------------------------------------

%Captital Letters
\newcommand{\mfA}{\mathfrak{A}}	\newcommand{\mfB}{\mathfrak{B}}
\newcommand{\mfC}{\mathfrak{C}}	\newcommand{\mfD}{\mathfrak{D}}
\newcommand{\mfE}{\mathfrak{E}}	\newcommand{\mfF}{\mathfrak{F}}
\newcommand{\mfG}{\mathfrak{G}}	\newcommand{\mfH}{\mathfrak{H}}
\newcommand{\mfI}{\mathfrak{I}}	\newcommand{\mfJ}{\mathfrak{J}}
\newcommand{\mfK}{\mathfrak{K}}	\newcommand{\mfL}{\mathfrak{L}}
\newcommand{\mfM}{\mathfrak{M}}	\newcommand{\mfN}{\mathfrak{N}}
\newcommand{\mfO}{\mathfrak{O}}	\newcommand{\mfP}{\mathfrak{P}}
\newcommand{\mfQ}{\mathfrak{Q}}	\newcommand{\mfR}{\mathfrak{R}}
\newcommand{\mfS}{\mathfrak{S}}	\newcommand{\mfT}{\mathfrak{T}}
\newcommand{\mfU}{\mathfrak{U}}	\newcommand{\mfV}{\mathfrak{V}}
\newcommand{\mfW}{\mathfrak{W}}	\newcommand{\mfX}{\mathfrak{X}}
\newcommand{\mfY}{\mathfrak{Y}}	\newcommand{\mfZ}{\mathfrak{Z}}
%Small Letters
\newcommand{\mfa}{\mathfrak{a}}	\newcommand{\mfb}{\mathfrak{b}}
\newcommand{\mfc}{\mathfrak{c}}	\newcommand{\mfd}{\mathfrak{d}}
\newcommand{\mfe}{\mathfrak{e}}	\newcommand{\mff}{\mathfrak{f}}
\newcommand{\mfg}{\mathfrak{g}}	\newcommand{\mfh}{\mathfrak{h}}
\newcommand{\mfi}{\mathfrak{i}}	\newcommand{\mfj}{\mathfrak{j}}
\newcommand{\mfk}{\mathfrak{k}}	\newcommand{\mfl}{\mathfrak{l}}
\newcommand{\mfm}{\mathfrak{m}}	\newcommand{\mfn}{\mathfrak{n}}
\newcommand{\mfo}{\mathfrak{o}}	\newcommand{\mfp}{\mathfrak{p}}
\newcommand{\mfq}{\mathfrak{q}}	\newcommand{\mfr}{\mathfrak{r}}
\newcommand{\mfs}{\mathfrak{s}}	\newcommand{\mft}{\mathfrak{t}}
\newcommand{\mfu}{\mathfrak{u}}	\newcommand{\mfv}{\mathfrak{v}}
\newcommand{\mfw}{\mathfrak{w}}	\newcommand{\mfx}{\mathfrak{x}}
\newcommand{\mfy}{\mathfrak{y}}	\newcommand{\mfz}{\mathfrak{z}}


\newcommand{\mytitle}{FMA\_3TC13\_TP Poly Exercices 2}
\newcommand{\myauthor}{github/omgeta}
\newcommand{\mydate}{AU 25-26 S1}

\begin{document}
\raggedright
\footnotesize
\begin{center}
{\normalsize{\textbf{\mytitle}}} \\
{\footnotesize{\mydate\hspace{2pt}\textemdash\hspace{2pt}\myauthor}}
\end{center}
\setlist{topsep=-1em, itemsep=-1em, parsep=2em}
\setlist[enumerate]{topsep=0.4em,itemsep=0.8em,parsep=0.35em}
%%%%%%%%%%%%%%%%%%%%%%%%%%%%%%%%%%%%%%%%%%%%%%%%%%%%%%
%                      Begin                         %
%%%%%%%%%%%%%%%%%%%%%%%%%%%%%%%%%%%%%%%%%%%%%%%%%%%%%%
\begin{enumerate}[label=2.\arabic*., start=17]
  \item \begin{minipage}[t]{\linewidth}\raggedright
\textbf{Question.} Le sous-ensemble $K=\{r+is:r,s\in\QQ\}\subseteq\CC$ est-il un corps pour l'addition et la multiplication usuelles ?
\par\smallskip
\textbf{Réponse.} $K$ est un corps pour les lois usuelles
    \begin{enumproof}
      \item Soient $z_1 = r_1 + is_1$ et $z_2 = r_2 + is_2$ dans $K$, où $r_1, s_1, r_2, s_2 \in \QQ$
      \item $z_1 + z_2 = (r_1 + r_2) + i(s_1 + s_2) \in K$\hfill(Stabilité par addition)
      \item $z_1z_2 = (r_1r_2 - s_1s_2) + i(r_1s_2 + s_1r_2) \in K$\hfill(Stabilité par multiplication)
      \item $0, 1 \in K$ et $-z_1 = -r_1 - is_1 \in K$
      \item Si $z_1 \neq 0$, $\displaystyle z_1^{-1} = \frac{r_1}{r_1^2 + s_1^2} - i\frac{s_1}{r_1^2 + s_1^2} \in K$
      \item Les autres propriétés sont héritées de $\CC$
      \item $\therefore K$ est un corps
    \end{enumproof}
\end{minipage}

  \item \begin{minipage}[t]{\linewidth}\raggedright
\textbf{Énoncé.} Soit $n\in\NN$. Établir les divisibilités suivantes.

\begin{enumerate}[\arabic*.]
        \item \textbf{Question.} $6\mid5n^3+n$.
\par\smallskip
\textbf{Réponse.} $5n^3 + n \equiv -n^3 + n = -n(n-1)(n+1) \pmod 6$

        Parmi trois entiers consécutifs, l'un est divisible par $3$ et l'un est pair.
        $\therefore 6 \mid 5n^3 + n$

  \item \textbf{Question.} $5\mid2^{2n+1}+3^{2n+1}$.
\par\smallskip
\textbf{Réponse.} $2^{2n+1} + 3^{2n+1} \equiv 2^{2n+1} + (-2)^{2n+1} = 0 \pmod 5$

        $\therefore 5 \mid 2^{2n+1} + 3^{2n+1}$

  \item \textbf{Question.} $11\mid2^{123}+3^{121}$.
\par\smallskip
\textbf{Réponse.} $2^{10} \equiv 1 \pmod{11}$ et $3^5 \equiv 1 \pmod{11}$, donc
        \[
          2^{123} + 3^{121} \equiv (2^{10})^{12}2^3 + (3^5)^{24}3
          \equiv 8 + 3 \equiv 0 \pmod{11}
        \]
        $\therefore 11 \mid 2^{123} + 3^{121}$

  \item \textbf{Question.} $9\mid4^n+6n-1$.
\par\smallskip
\textbf{Réponse.} Par le binôme de Newton, $4^n = (1+3)^n \equiv 1 + 3n \pmod 9$, car les termes de degré au moins $2$ en $3$ sont divisibles par $9$.

        Alors $4^n + 6n - 1 \equiv 1 + 3n + 6n - 1 = 9n \equiv 0 \pmod 9$

        $\therefore 9 \mid 4^n + 6n - 1$

\end{enumerate}
\end{minipage}

  \item \begin{minipage}[t]{\linewidth}\raggedright
\textbf{Énoncé.} Résoudre dans $\ZZ^2$ les équations suivantes.

\begin{enumerate}[\arabic*.]
        \item \textbf{Question.} $xy=3x+2y$.
\par\smallskip
\textbf{Réponse.} $xy = 3x + 2y \iff (x-2)(y-3) = 6$

        On pose $x-2 = d$ et $y-3 = 6/d$, où $d \in \{\pm 1, \pm 2, \pm 3, \pm 6\}$.

        $\therefore (x, y) \in \{(3,9), (4,6), (5,5), (8,4), (1,-3), (0,0), (-1,1), (-4,2)\}$

  \item \textbf{Question.} $\displaystyle\frac1x+\frac1y=\frac15$.
\par\smallskip
\textbf{Réponse.} Pour $x, y \neq 0$, $\displaystyle \frac{1}{x} + \frac{1}{y} = \frac{1}{5} \iff (x-5)(y-5) = 25$

        On pose $x-5 = d$ et $y-5 = 25/d$, où $d \in \{\pm 1, \pm 5, \pm 25\}$.
        Le cas $d = -5$ donne $(0,0)$, qui est exclu.

        $\therefore (x, y) \in \{(6,30), (10,10), (30,6), (4,-20), (-20,4)\}$

  \item \textbf{Question.} $x^2-y^2-4x-2y-5=0$.
\par\smallskip
\textbf{Réponse.} $x^2 - y^2 - 4x - 2y - 5 = 0 \iff (x-2)^2 - (y+1)^2 = 8$
        \[
          \iff (x-y-3)(x+y-1) = 8
        \]
        On pose $x-y-3 = d$ et $x+y-1 = 8/d$, où $d \mid 8$. Alors
        \[
          x = 2 + \frac{d + 8/d}{2}, \qquad y = -1 + \frac{8/d - d}{2}
        \]
        Pour que $x, y \in \ZZ$, $d$ et $8/d$ doivent avoir la même parité, donc $d \in \{\pm 2, \pm 4\}$.

        $\therefore (x, y) \in \{(5,0), (5,-2), (-1,0), (-1,-2)\}$

\end{enumerate}
\end{minipage}

  \item \begin{minipage}[t]{\linewidth}\raggedright
\begin{enumerate}[\arabic*.]
        \item \textbf{Question.} Peut-on trouver un inverse de $117$ modulo $368$ ?
\par\smallskip
\textbf{Réponse.} $117 \cdot 173 - 55 \cdot 368 = 1$\hfill(Relation de Bézout)

        Alors $117 \cdot 173 \equiv 1 \pmod{368}$.
        $\therefore 117^{-1} \equiv 173 \pmod{368}$

  \item \textbf{Question.} $\ZZ/368\ZZ$ est-il un corps ?
\par\smallskip
\textbf{Réponse.} $368 = 2^4 \cdot 23$ n'est pas premier.
        $\therefore \ZZ/368\ZZ$ n'est pas un corps

\end{enumerate}
\end{minipage}

  \item \begin{minipage}[t]{\linewidth}\raggedright
\textbf{Question.} Dresser la table des inverses dans le corps $\ZZ/23\ZZ$.
\par\smallskip
\textbf{Réponse.} Dans $\ZZ/23\ZZ$, $0$ n'a pas d'inverse et les autres inverses sont :
    \[
      \begin{array}{c|rrrrrrrrrrr}
        x & 1 & 2 & 3 & 4 & 5 & 6 & 7 & 8 & 9 & 10 & 11\\ \hline
        x^{-1} & 1 & 12 & 8 & 6 & 14 & 4 & 10 & 3 & 18 & 7 & 21
      \end{array}
    \]
    \[
      \begin{array}{c|rrrrrrrrrrr}
        x & 12 & 13 & 14 & 15 & 16 & 17 & 18 & 19 & 20 & 21 & 22\\ \hline
        x^{-1} & 2 & 16 & 5 & 20 & 13 & 19 & 9 & 17 & 15 & 11 & 22
      \end{array}
    \]
\end{minipage}

  \item \begin{minipage}[t]{\linewidth}\raggedright
\textbf{Question.} Résoudre $9x+3\equiv0\pmod{22}$.
\par\smallskip
\textbf{Réponse.} $5 \cdot 9 - 2 \cdot 22 = 1$, donc $9^{-1} \equiv 5 \pmod{22}$.
    \[
      9x + 3 \equiv 0 \pmod{22} \iff x \equiv -3 \cdot 5 \equiv 7 \pmod{22}
    \]
    $\therefore x = 7 + 22k$, où $k \in \ZZ$
\end{minipage}

  \item \begin{minipage}[t]{\linewidth}\raggedright
\textbf{Question.} Déterminer $\operatorname{pgcd}(2^{445}+7,15)$.
\par\smallskip
\textbf{Réponse.} $2^4 \equiv 1 \pmod{15}$ et $445 = 4 \cdot 111 + 1$, donc
    \[
      2^{445} + 7 \equiv (2^4)^{111}2 + 7 \equiv 9 \pmod{15}
    \]
    $\therefore \operatorname{pgcd}(2^{445} + 7, 15) = \operatorname{pgcd}(9, 15) = 3$
\end{minipage}

  \item \begin{minipage}[t]{\linewidth}\raggedright
\textbf{Question.} Soit $n\geq1$. Montrer que $\operatorname{pgcd}(2n+4,3n+3)\in\{1,2,3,6\}$. Les quatre cas sont-ils possibles ?
\par\smallskip
\textbf{Réponse.} Soit $d = \operatorname{pgcd}(2n+4, 3n+3)$
    \begin{enumproof}
      \item $3(2n+4) - 2(3n+3) = 6$
      \item $d \mid 6$\hfill(Combinaison linéaire)
      \item $\therefore d \in \{1, 2, 3, 6\}$
    \end{enumproof}
    Les quatre cas sont possibles :
    \[
      \begin{array}{c|c|c|c}
        n & 2n+4 & 3n+3 & \operatorname{pgcd}\\ \hline
        1 & 6 & 6 & 6\\
        2 & 8 & 9 & 1\\
        3 & 10 & 12 & 2\\
        4 & 12 & 15 & 3
      \end{array}
    \]
\end{minipage}

  \item \begin{minipage}[t]{\linewidth}\raggedright
\textbf{Question.} Soit $n\geq1$. Montrer que $3n^2+2n$ et $n+1$ sont premiers entre eux.
\par\smallskip
\textbf{Réponse.} $a - (3n-1)b = 3n^2 + 2n - (3n-1)(n+1) = 1$

    Tout diviseur commun de $a$ et $b$ divise $1$.
    $\therefore \operatorname{pgcd}(a,b) = 1$
\end{minipage}

  \item \begin{minipage}[t]{\linewidth}\raggedright
\textbf{Question.} Soit $n\geq1$. Montrer que $n^2+n$ et $2n+1$ sont premiers entre eux.
\par\smallskip
\textbf{Réponse.} $-4a + (2n+1)b = -4(n^2+n) + (2n+1)^2 = 1$

    $\therefore \operatorname{pgcd}(a,b) = 1$\hfill(Relation de Bézout)
\end{minipage}

  \item \begin{minipage}[t]{\linewidth}\raggedright
\textbf{Question.} Soient $n\geq1$ et $\alpha\neq\beta$ dans $\{1,\ldots,n+1\}$. Montrer que $\alpha n!+1$ et $\beta n!+1$ sont premiers entre eux.
\par\smallskip
\textbf{Réponse.} $a = \alpha n! + 1$ et $b = \beta n! + 1$ sont premiers entre eux
    \begin{enumproof}
      \item Supposons qu'un nombre premier $p$ divise $a$ et $b$
      \item $p \mid a-b = (\alpha-\beta)n!$
      \item Alors $p \mid (\alpha-\beta)$ ou $p \mid n!$\hfill($p$ est premier)
      \item Comme $1 \leq |\alpha-\beta| \leq n$, on a $p \leq n$ dans les deux cas
      \item Donc $p \mid n!$ et $a = \alpha n! + 1 \equiv 1 \pmod p$
      \item Ceci contredit $p \mid a$
      \item $\therefore \operatorname{pgcd}(a,b) = 1$
    \end{enumproof}
\end{minipage}

  \item \begin{minipage}[t]{\linewidth}\raggedright
\begin{enumerate}[\arabic*.]
        \item \textbf{Question.} Tous les éléments de $\ZZ/8\ZZ$ sont-ils inversibles pour la multiplication ?
\par\smallskip
\textbf{Réponse.} Non, par exemple $2$ n'est pas inversible car $\operatorname{pgcd}(2,8) = 2 \neq 1$.

  \item \textbf{Question.} $\ZZ/8\ZZ$ est-il un corps ?
\par\smallskip
\textbf{Réponse.} $\ZZ/8\ZZ$ n'est pas un corps car $8$ n'est pas premier.

  \item \textbf{Question.} Donner les inverses de tous les éléments inversibles de $\ZZ/8\ZZ$.
\par\smallskip
\textbf{Réponse.} Les éléments inversibles sont $1, 3, 5, 7$, chacun étant son propre inverse :
        \[
          1^2 \equiv 3^2 \equiv 5^2 \equiv 7^2 \equiv 1 \pmod 8
        \]

\end{enumerate}
\end{minipage}

  \item \begin{minipage}[t]{\linewidth}\raggedright
\textbf{Énoncé.} On munit $A=\{\alpha+\beta\sqrt2:\alpha,\beta\in\ZZ\}\subseteq\RR$ des lois usuelles et on pose $N(\alpha+\beta\sqrt2)=(\alpha+\beta\sqrt2)(\alpha-\beta\sqrt2)$.

\begin{enumerate}[\arabic*.]
        \item \textbf{Question.} Montrer que $N(zz^{\prime})=N(z)N(z^{\prime})$ pour tous $z,z^{\prime}\in A$.
\par\smallskip
\textbf{Réponse.} Pour $z = \alpha + \beta\sqrt 2 \in A$, on note $\overline z = \alpha - \beta\sqrt 2$, donc $N(z) = z\overline z = \alpha^2 - 2\beta^2 \in \ZZ$.

Soient $z = \alpha + \beta\sqrt 2$ et $z' = \alpha' + \beta'\sqrt 2$ dans $A$
        \begin{enumproof}
          \item $zz' = (\alpha\alpha' + 2\beta\beta') + (\alpha\beta' + \beta\alpha')\sqrt 2 \in A$
          \item $\overline{zz'} = (\alpha\alpha' + 2\beta\beta') - (\alpha\beta' + \beta\alpha')\sqrt 2 = \overline z\,\overline{z'}$
          \item $\therefore N(zz') = zz'\overline{zz'} = (z\overline z)(z'\overline{z'}) = N(z)N(z')$
        \end{enumproof}

  \item \textbf{Question.} Calculer les inverses de $3+2\sqrt2$ et $7+5\sqrt2$, et vérifier leur appartenance à $A$.
\par\smallskip
\textbf{Réponse.} $\displaystyle (3+2\sqrt 2)^{-1} = \frac{3-2\sqrt 2}{9-8} = 3-2\sqrt 2 \in A$

        $\displaystyle (7+5\sqrt 2)^{-1} = \frac{7-5\sqrt 2}{49-50} = -7+5\sqrt 2 \in A$

  \item \textbf{Question.} Montrer que $z\in A$ est inversible dans $A$ si et seulement si $N(z)\in\{1,-1\}$.
\par\smallskip
\textbf{Réponse.} $z$ est inversible dans $A$ si et seulement si $N(z) \in \{1,-1\}$
        \begin{enumproof}
          \item ($\Rightarrow$) Si $zw = 1$ avec $w \in A$, alors $N(z)N(w) = N(1) = 1$.
            Comme ces normes sont entières, $N(z) \in \{1,-1\}$.
          \item ($\Leftarrow$) Si $N(z) \in \{1,-1\}$, alors $\displaystyle z^{-1} = \frac{\overline z}{N(z)} = \pm\overline z \in A$.
        \end{enumproof}

  \item \textbf{Question.} $A$ est-il un corps pour les lois usuelles ?
\par\smallskip
\textbf{Réponse.} $N(2) = 4 \notin \{1,-1\}$, donc $2 \neq 0$ n'est pas inversible dans $A$.

        $\therefore A$ n'est pas un corps

\end{enumerate}
\end{minipage}

  \item \begin{minipage}[t]{\linewidth}\raggedright
\textbf{Énoncé.} Soit $p$ un nombre premier.

\begin{enumerate}[\arabic*.]
        \item \textbf{Question.} Pour $1\leq k\leq p-1$, montrer que $p\mid\binom pk$.
\par\smallskip
\textbf{Réponse.} Soit $k \in \{1,\ldots,p-1\}$
        \begin{enumproof}
          \item $\displaystyle k!(p-k)!\binom pk = p!$
          \item $p \nmid k!(p-k)!$ car tous les facteurs sont strictement inférieurs à $p$
          \item Comme $p$ est premier et $p \mid p!$, $\displaystyle p \mid \binom pk$
        \end{enumproof}

  \item \textbf{Question.} Pour $a,b\in\ZZ/p\ZZ$, montrer que $(a+b)^p\equiv a^p+b^p\pmod p$.
\par\smallskip
\textbf{Réponse.} Par le binôme de Newton et la question 1,
        \[
          (a+b)^p = \sum_{k=0}^p \binom pk a^kb^{p-k}
          \equiv a^p + b^p \pmod p
        \]

  \item \textbf{Question.} Montrer que $a^p\equiv a\pmod p$ pour tout $a\in\ZZ/p\ZZ$.
\par\smallskip
\textbf{Réponse.} Par récurrence sur un représentant $a \in \{0,\ldots,p-1\}$
        \begin{enumproof}
          \item $0^p \equiv 0 \pmod p$\hfill(Initialisation)
          \item Si $a^p \equiv a \pmod p$, alors $(a+1)^p \equiv a^p + 1^p \equiv a+1 \pmod p$\hfill(Question 2)
          \item $\therefore a^p \equiv a \pmod p$ pour tout $a \in \ZZ/p\ZZ$
        \end{enumproof}

\end{enumerate}
\end{minipage}

  \item \begin{minipage}[t]{\linewidth}\raggedright
\textbf{Énoncé.} Soient $a,b,u,v\in\ZZ$.

\begin{enumerate}[\arabic*.]
        \item \textbf{Question.} Vérifier $(au+bv)^2=(a+b)(au^2+bv^2)+ab(2uv-u^2-v^2)$.
\par\smallskip
\textbf{Réponse.} En développant le membre de droite,
        \begin{align*}
          (a+b)(au^2+bv^2) + ab(2uv-u^2-v^2)
          &= a^2u^2 + b^2v^2 + 2abuv\\
          &= (au+bv)^2
        \end{align*}

  \item \textbf{Question.} En déduire que si $a,b$ sont premiers entre eux, alors $a+b,ab$ le sont aussi.
\par\smallskip
\textbf{Réponse.} Si $\operatorname{pgcd}(a,b) = 1$, alors $\operatorname{pgcd}(a+b,ab) = 1$
        \begin{enumproof}
          \item Il existe $u,v \in \ZZ$ tels que $au+bv = 1$\hfill(Bézout)
          \item Soit $d \geq 1$ un diviseur commun de $a+b$ et $ab$
          \item $d \mid (a+b)(au^2+bv^2) + ab(2uv-u^2-v^2) = (au+bv)^2 = 1$
          \item $\therefore d = 1$ et $\operatorname{pgcd}(a+b,ab) = 1$
        \end{enumproof}

  \item \textbf{Question.} Montrer la réciproque : si $a+b,ab$ sont premiers entre eux, alors $a,b$ le sont aussi.
\par\smallskip
\textbf{Réponse.} Si $\operatorname{pgcd}(a+b,ab) = 1$, alors $\operatorname{pgcd}(a,b) = 1$
        \begin{enumproof}
          \item Soit $d \geq 1$ un diviseur commun de $a$ et $b$
          \item $d \mid a+b$ et $d \mid ab$, donc $d \mid \operatorname{pgcd}(a+b,ab) = 1$
          \item $\therefore d = 1$ et $\operatorname{pgcd}(a,b) = 1$
        \end{enumproof}

\end{enumerate}
\end{minipage}

  \item \begin{minipage}[t]{\linewidth}\raggedright
\begin{enumerate}[\arabic*.]
        \item \textbf{Question.} Montrer que $\sqrt2\notin\QQ$.
\par\smallskip
\textbf{Réponse.} $\sqrt 2 \notin \QQ$
        \begin{enumproof}
          \item Supposons $\sqrt 2 = r/s$, où $r,s \in \NN$, $s \neq 0$ et $\operatorname{pgcd}(r,s) = 1$
          \item $r^2 = 2s^2$, donc $r$ est pair et s'écrit $r = 2t$
          \item Alors $4t^2 = 2s^2$, donc $s^2 = 2t^2$ et $s$ est pair
          \item Ceci contredit $\operatorname{pgcd}(r,s) = 1$
          \item $\therefore \sqrt 2 \notin \QQ$
        \end{enumproof}

  \item \textbf{Question.} Pour tout $n\in\NN$, montrer que $(1+\sqrt2)^n=a_n+b_n\sqrt2$ pour un unique couple $(a_n,b_n)\in\NN^2$.
\par\smallskip
\textbf{Réponse.} Existence et unicité de $(a_n,b_n) \in \NN^2$
        \begin{enumproof}
          \item Pour $n = 0$, $(1+\sqrt 2)^0 = 1$, donc $(a_0,b_0) = (1,0)$
          \item Si $(1+\sqrt 2)^n = a_n + b_n\sqrt 2$, alors
            \[
              (1+\sqrt 2)^{n+1} = (a_n+2b_n) + (a_n+b_n)\sqrt 2
            \]
            Donc $a_{n+1} = a_n+2b_n$ et $b_{n+1} = a_n+b_n$ sont dans $\NN$.
          \item Par récurrence, le couple existe pour tout $n \in \NN$
          \item Si $a+b\sqrt 2 = a'+b'\sqrt 2$ et $b \neq b'$, alors $\displaystyle \sqrt 2 = \frac{a'-a}{b-b'} \in \QQ$, contradiction
          \item Donc $b = b'$ puis $a = a'$. $\therefore$ le couple est unique
        \end{enumproof}

  \item \textbf{Question.} Calculer $a_n^2-2b_n^2$.
\par\smallskip
\textbf{Réponse.} La conjugaison donne $(1-\sqrt 2)^n = a_n-b_n\sqrt 2$, donc
        \[
          a_n^2 - 2b_n^2 = (a_n+b_n\sqrt 2)(a_n-b_n\sqrt 2)
          = \big((1+\sqrt 2)(1-\sqrt 2)\big)^n = (-1)^n
        \]

  \item \textbf{Question.} En déduire que $a_n,b_n$ sont premiers entre eux.
\par\smallskip
\textbf{Réponse.} $(-1)^n a_n^2 - 2(-1)^n b_n^2 = 1$\hfill(Question 3)

        C'est une relation de Bézout entre $a_n$ et $b_n$.
        $\therefore \operatorname{pgcd}(a_n,b_n) = 1$

\end{enumerate}
\end{minipage}

\end{enumerate}
%%%%%%%%%%%%%%%%%%%%%%%%%%%%%%%%%%%%%%%%%%%%%%%%%%%%%%
%                       End                          %
%%%%%%%%%%%%%%%%%%%%%%%%%%%%%%%%%%%%%%%%%%%%%%%%%%%%%%

\end{document}
